% !TeX root = ../main.tex
\section{结论}
\label{sec:conclusion}

DP-Puzzle 通过有效数重编码与分级对齐共享乘法、移位和累加，并以 FP4 专用化及 $K/N$ 组合扩展低精度吞吐。共享降低面积，专用化抑制宽通路开销，两者共同决定增量算力的效率。固定面积分析表明，配比选择须同时考虑核心数量和各精度的单位计算能耗，中间峰值配比未必具有独立优势。三个 DeepSeek 模型的参考 W4A8 prefill 均受益于 $1:4:8$；采用 NVFP4 W4A4 路由专家后，V4-Flash 的长上下文负载已转向 $1:2:8$，其余模型仍偏向 $1:4:8$。配比应随实际执行精度演进，匹配各阶段对时间和能耗的贡献。
